In Louscoone Inlet,
local herring stocks would normally be found well up inside the inlet and
would often present a very light spawn near the head of the inlet early
in the season (early- to mid-March).
This was the case in 2025.
Herring were sounded in Louscoone Inlet on March \nth{12} and \nth{13} in small dispersed schools.
There was a spawn reported on March 12 toward the head of the inlet
along the East shore which was believed to be very light.

The majority of herring stocks in the Haida Gwaii major SAR
continue to be located around Burnaby Island (in Skincuttle Inlet and upper Burnaby Strait).
The area with the most abundant spawn in 2025 was in Skaat Harbour and out toward Newberry Cove.
The abundance of stocks and deposition of spawns appear to be very similar to past years (2022 to 2024).
These stocks appear to be maintaining themselves but are not really showing any signs
of meaningful growth despite no active fisheries for over 20 years.

The Juan Perez Sound inlets normally support spawns when
there are a lot of `young of the year' fish in the population.
There were no spawns observed in the upper Juan Perez Sound inlets (North of Werner Point)
during the 2025 season.

In most years Atli Inlet is found to support a small stock of herring.
Such was the case in 2025 with 500 tons sounded and sampled on April \nth{3} and
a light spawn in Beljay Bay which occurred mid-April. 

Herring stocks in Selwyn Inlet continue to show declining trends.
In recent years (2021 and 2023) the spawns in Selwyn Inlet occurred
late in the season (end of April) and were light in intensity.
In 2024 spawns were almost non-existant
(only a very light spot spawn near Traynor Creek mid-April).
In 2025 there were spawns observed/reported April \nth{8} and \nth{9}
in Pacofi Bay and Cecil Cove, which was better than in recent years,
but still minimal in comparison to historic Selwyn Inlet spawns
(adding to concerns that the Selwyn Inlet herring stocks
may be in a serious and irreversible decline).

Herring spawns in Cumshewa Inlet have been infrequent in recent years.
During the 2025 season there was a spawn observed/reported early in the season (March \nth{25})
toward the outer portion of the inlet along the North shore and
North of Cumshewa Head toward Grey Bay.

Approximately 5,000 kg of k'aaw was harvested from the major SAR
of East Moresby (primarily upper Burnaby Strait) in early- to mid-April.
Product was reported to be moderate to very good in quality.
Most of the traditional k'aaw was harvested by two Haida individuals, however,
much of the product was distributed throughout the communities of Skidegate and Old Massett.
There have been no known harvests of k'aaw from other locations on Haida Gwaii
(i.e. Selwyn Inlet or Skidegate Inlet).
